\chapter*{Remerciements}

Au terme de ce stage, il m’est particulièrement agréable d’exprimer ma profonde gratitude à toutes celles et à tous ceux qui, de près ou de loin, ont contribué à la bonne réalisation de ce rapport.\\

Je tiens à exprimer mes sincères remerciements au Dr Abdou \textsc{DIOUF}, Directeur Général de l’ANSD, ainsi qu’au Dr Momath \textsc{CISSÉ},Directeur Général Adjoint de l’ANSD et à M. Fodé \textsc{DIÈME}, directeur de la DMCI pour m’avoir fait l’honneur de m’accueillir au sein de cette institution.\\

J’adresse également mes remerciements à M. \textsc{NDIASSE}, Chef de la division de la Méthodologie et de l’Innovation, ainsi qu’à M. \textsc{AAAAAA}, Responsable du Bureau des Méthodes et Normes (BMN), structure au sein de laquelle j’ai eu le privilège d’effectuer mon stage.\\

Mes remerciements les plus chaleureux vont également à mes encadreurs, M. Mamadou \textsc{GUEYE}, M. Ouseynou \textsc{MBAYE}, M. Cheikh \textsc{GUEYE} Ingénieurs statisticiens économistes à la DMCI, pour leur disponibilité, leurs conseils avisés et le suivi constant dont j’ai bénéficié tout au long de ces deux mois.\\  

Je souhaite aussi exprimer ma reconnaissance à Mme \textsc{Ndeye}, M. Libasse \textsc{-----},, ainsi qu’à l’ensemble du personnel de la DMCI pour leur accueil chaleureux, leurs échanges enrichissants et leur ouverture d’esprit.\\

Je ne saurais oublier de remercier l’ensemble du corps administratif et professoral de l’ENSAE, en particulier M. Idrissa \textsc{DIAGNE}, Directeur de l’ENSAE, Dr Alassane \textsc{AW}, Responsable de la filière des Analystes Statisticiens (AS), M. Mamadou \textsc{BALDÉ}, Responsable de la filière des Ingénieurs Statisticiens Économistes (ISE)%, et Dr Souleymane \textsc{FOFANA}, Chef de l’Unité de recherche.\\

Enfin, j’adresse mes remerciements les plus chaleureux à mes camarades de la promotion \textsc{AS 2023--2026}, pour l’esprit de solidarité et de fraternité qui nous unit, ainsi qu’à toutes les personnes qui, bien que non citées nommément, ont apporté leur soutien et contribué, d’une manière ou d’une autre, à l’aboutissement de ce travail.  

